\documentclass[12pt,                       
               twoside,
               a4paper,                   
               BCOR=1.0cm,                 
               headsepline,
               openany,              
               toc = index,            
               bibliography=totoc,          
               numbers=noenddot,            
               ]{scrbook}
\usepackage[utf8]{inputenc} 
				
\usepackage{scrhack}					  
\frenchspacing     
\usepackage{graphicx}
\usepackage[margin=10pt,font=small,labelfont=bf,justification=centering]{caption}
\usepackage{xcolor}	
% Define a command for a blue section heading
\newcommand{\bluesec}[1]{\section{\textcolor{blue}{#1}}}	
%\usepackage{subfigure}
\usepackage{subcaption}
\usepackage[export]{adjustbox}
\usepackage{float}    
\usepackage{array}    
\usepackage{hhline}   
\usepackage{tabularx} 
\usepackage{longtable}
\usepackage[intlimits,sumlimits]{amsmath} 
\usepackage{amssymb}
\usepackage{makeidx}
\usepackage{showidx}
\usepackage{enumerate}
\usepackage{multirow}           
\usepackage[colorlinks,pdfpagelabels,pdfstartview = FitH,bookmarksopen = true,bookmarksnumbered = true,linkcolor = black,plainpages = false,hypertexnames = false,citecolor = black]{hyperref}
\usepackage{esvect}
\usepackage{amsthm}
\usepackage{amsfonts}					
\usepackage{textgreek}                
\usepackage{siunitx}           
\usepackage{xspace}
\usepackage{ragged2e}
\makeindex

\pdfsuppresswarningpagegroup=1
                                
\PassOptionsToPackage{american}{babel}  
\usepackage{babel}

\usepackage{listings}
%\lstset{emph={trueIndex,root},emphstyle=\color{BlueViolet}}%\underbar} % for special keywords
\lstset{language=[LaTeX]Tex,%C++,
	morekeywords={PassOptionsToPackage,selectlanguage},
	keywordstyle=\color{RoyalBlue},%\bfseries,
	basicstyle=\small\ttfamily,
	%identifierstyle=\color{NavyBlue},
	commentstyle=\color{Green}\ttfamily,
	stringstyle=\rmfamily,
	numbers=none,%left,%
	numberstyle=\scriptsize,%\tiny
	stepnumber=5,
	numbersep=8pt,
	showstringspaces=false,
	breaklines=true,
	%frameround=ftff,
	%frame=single,
	belowcaptionskip=.75\baselineskip
	%frame=L
}

% VHDL code formatting settings
\lstdefinelanguage{VHDL}{
	keywords=[1]{architecture, signal, component, process, begin, case, when, others, end, entity, port, if, is, else},
	keywords=[2]{std\_logic, std\_logic\_vector, unsigned, signed, clk},
	sensitive=false,
	morecomment=[l]--,       % Single-line comments
	morestring=[b]",         % Strings in double quotes
}

% Global lstlisting settings for code formatting
\lstset{
	language=VHDL,                   % Set VHDL as the default language for listings
	basicstyle=\ttfamily\footnotesize,% Code font style and size
	keywordstyle=\color{blue}\bfseries, % Keywords in blue and bold
	commentstyle=\color{gray},       % Comments in gray
	stringstyle=\color{orange},      % Strings in orange
	numbers=left,                    % Line numbers on the left
	numberstyle=\tiny\color{gray},   % Line number style
	stepnumber=1,                    % Increment line numbers by 1
	frame=single,                    % Frame around the code
	breaklines=true,                 % Automatically break long lines
	tabsize=4,                       % Set tab size to 4 spaces
	captionpos=b                     % Caption at the bottom of listings
}


\usepackage[printonlyused]{acronym}
\usepackage{csquotes}
\usepackage{gensymb}

% Adjust paragraph spacing and indentation
\setlength{\parindent}{0pt} % No paragraph indentation
\setlength{\parskip}{0.5\baselineskip}   % No extra vertical space between paragraphs


					
